\exercisesheader{}

% 23

\eoce{\qt{Kelereng di dalam guci\label{marbles_in_urn}} Bayangkan Anda memiliki sebuah guci 
yang berisi 5 kelereng merah, 3 biru, dan 2 oranye. 
\begin{parts}
\item Berapakah peluang bahwa kelereng pertama yang Anda ambil berwarna biru?
\item Misalkan Anda mengambil kelereng biru pada pengambilan pertama. Jika pengambilan dilakukan dengan 
pengembalian, berapakah peluang mengambil kelereng biru pada pengambilan kedua?
\item Misalkan Anda justru mengambil kelereng oranye pada pengambilan pertama. Jika pengambilan dilakukan 
dengan pengembalian, berapakah peluang mengambil kelereng biru pada pengambilan 
kedua?
\item Jika pengambilan dilakukan dengan pengembalian, berapakah peluang mengambil dua kelereng 
biru berturut-turut?
\item Ketika pengambilan dilakukan dengan pengembalian, apakah setiap pengambilan saling independen? Jelaskan.
\end{parts}
}{}

% 24

\eoce{\qt{Kaus kaki di dalam laci\label{socks_in_drawer}} Di dalam laci kaus kaki, Anda memiliki 
4 kaus kaki biru, 5 abu-abu, dan 3 hitam. Dalam keadaan setengah mengantuk pada suatu pagi, Anda mengambil 2 kaus kaki secara 
acak lalu memakainya. Tentukan peluang bahwa Anda akhirnya mengenakan
\begin{parts}
\item 2 kaus kaki biru
\item tidak ada kaus kaki abu-abu
\item setidaknya 1 kaus kaki hitam
\item sebuah kaus kaki hijau
\item sepasang kaus kaki dengan warna yang sama
\end{parts}
}{}

% 25

\eoce{\qt{Keping di dalam kantong\label{chips_in_bag}} Bayangkan Anda memiliki sebuah kantong 
yang berisi 5 keping merah, 3 biru, dan 2 oranye.
\begin{parts}
\item Misalkan Anda mengambil satu keping dan keping itu berwarna biru. Jika pengambilan dilakukan tanpa pengembalian, 
berapakah peluang bahwa keping berikutnya juga berwarna biru?
\item Misalkan Anda mengambil satu keping dan keping itu berwarna oranye, lalu Anda mengambil keping kedua 
tanpa pengembalian. Berapakah peluang bahwa keping kedua ini berwarna biru?
\item Jika pengambilan dilakukan tanpa pengembalian, berapakah peluang mengambil dua keping biru 
berturut-turut?
\item Ketika pengambilan dilakukan tanpa pengembalian, apakah setiap pengambilan saling independen? Jelaskan.
\end{parts}
}{}

% 26

\eoce{\qt{Buku di rak buku\label{books_on_shelf}} Tabel di bawah ini menunjukkan 
distribusi buku di sebuah rak berdasarkan jenis fiksi atau 
nonfiksi serta format sampul keras atau sampul lunak.
\begin{center}
\begin{tabular}{ll  cc c} 
                                &           & \multicolumn{2}{c}{\textit{Format}} \\
\cline{3-4}
                                &           & Sampul keras     & Sampul lunak     & Total \\
\cline{2-5}
\multirow{2}{*}{\textit{Jenis}}  & Fiksi   & 13            & 59            & 72 \\
                                & Nonfiksi& 15            & 8             & 23 \\
\cline{2-5} 
                                & Total     & 28            & 67            & 95 \\
\cline{2-5}
\end{tabular}
\end{center}
\begin{parts}
\item Tentukan peluang mengambil sebuah buku sampul keras terlebih dahulu, lalu sebuah buku fiksi 
bersampul lunak sebagai buku kedua, ketika pengambilan dilakukan tanpa pengembalian.
\item Tentukan peluang mengambil sebuah buku fiksi terlebih dahulu, lalu sebuah 
buku sampul keras sebagai buku kedua, ketika pengambilan dilakukan tanpa pengembalian.
\item Hitung peluang untuk skenario pada bagian~(b), tetapi kali ini 
lakukan perhitungan untuk skenario ketika buku pertama dikembalikan 
ke rak sebelum buku kedua diambil secara acak.
\item Jawaban akhir untuk bagian~(b) dan~(c) sangat mirip. Jelaskan mengapa 
hal ini terjadi.
\end{parts}
}{}

% 27

\eoce{\qt{Pakaian mahasiswa\label{student_outfits}} Di sebuah kelas yang terdiri atas 24 
mahasiswa, 7 mahasiswa mengenakan jins, 4 mengenakan celana pendek, 8 mengenakan 
rok, dan sisanya mengenakan legging. Jika kita memilih 3 mahasiswa secara acak 
tanpa pengembalian, berapakah peluang bahwa satu mahasiswa yang terpilih 
mengenakan legging dan dua lainnya mengenakan jins? Perhatikan bahwa jenis-jenis 
pakaian ini saling lepas.
}{}

% 28

\eoce{\qt{Masalah ulang tahun\label{birthday_problem}} Misalkan kita memilih tiga 
orang secara acak. Untuk setiap pertanyaan berikut, abaikan kasus khusus 
ketika seseorang mungkin lahir pada 29 Februari, dan anggap kelahiran tersebar 
merata sepanjang tahun.
\begin{parts}
\item Berapakah peluang bahwa dua orang pertama berulang tahun pada tanggal yang sama? 
\item Berapakah peluang bahwa setidaknya dua orang berulang tahun pada tanggal yang sama?
\end{parts}
}{}
